\section{Arena Execution and Reproducibility}
\label{sec:arena}

The Arena platform implements the benchmark contract across participant interaction, market-data access, strategy execution, and final scoring. Its control plane records suites, submissions, replicates, operations, resource usage, and results. Quantitative workers and data services execute the authorized computation. A batch campaign interface uses these same services to compare declared combinations of models, harnesses, budgets, and suites.

\subsection{Separation of Research and Quantitative Execution}
\label{sec:arena_execution}

The research harness controls what to propose and when to evaluate or finish. It receives model and research tools, while computational strategies execute under a separate identity with scoped market-data access. A trusted evaluator consumes the resulting scores and applies the frozen financial rules. This division allows an agent to conduct data-intensive research through executable programs without forwarding raw market matrices or private evaluation artifacts into the model conversation. Table~\ref{tab:arena_roles} summarizes the responsibilities.

\begin{table}[t]
    \centering
    \small
    \caption{Responsibilities within the Arena. The final choice belongs to the participant; execution permissions, accounting, and assessment belong to the platform.}
    \label{tab:arena_roles}
    \begin{tabular}{@{}p{0.20\linewidth}p{0.35\linewidth}p{0.36\linewidth}@{}}
        \toprule
        Component & Controlled decisions & Recorded or returned evidence \\
        \midrule
        Research harness & Prompts, candidate revisions, stopping, final selection & Model interactions, research actions, selected strategy \\
        Strategy execution & Computation within the authorized input and runtime contract & Daily scores, declared features, optional fitted artifact \\
        Model gateway & Model access, operation identity, usage settlement & Provider response, observed usage, unresolved-operation state \\
        Data and execution services & Version binding, date and field scope, process lifecycle & Authorized execution records, resource observations \\
        Trusted evaluator & Labels, trading rules, metrics, qualification & Core results, coverage, private detailed evidence \\
        Campaign interface & Registered comparison matrix and collection & Existing scorecards, complete replicate and failure registry \\
        \bottomrule
    \end{tabular}
\end{table}

Model-provider credentials and data-access capabilities are held by trusted services. In the reusable execution environment, a supervisor advances the frozen sequence of signal dates. Each date starts a fresh low-privilege strategy process, working directory, and date-scoped data gateway. Before advancing, the supervisor closes the day's processes and connections and persists its accepted output. Reusing the virtual machine therefore preserves controlled execution infrastructure; participant state crosses stages only through declared artifacts such as a frozen fitted model. Earlier per-day runtime contracts remain identifiable by their registered versions.

Training, research, and final assessment use separately bound executions. The final evaluator accepts registered sealed releases through a private data plane, while research credentials retain their authorized research scope. The resulting boundary combines service authorization with process isolation. Claims about a hosted campaign additionally depend on the deployed runtime, network configuration, and boundary checks recorded for that campaign.

\subsection{Immutable Data and Shared Financial Evaluation}
\label{sec:arena_data}

Market observations are served from immutable publications. Offline preparation establishes the data layout, indexes, calendar, and declared coverage. Online access resolves the authorized release and reads bounded dates, instruments, and fields. Daily data use published queryable stores; indexed intraday data support column and time-range selection. Exact price and quantity semantics, exchange-local timestamps, and source missingness are retained through the query interface. A change to the default publication does not redirect an active suite away from its registered version.

Arena and the internal research platform use a shared daily replay and statistical core. Within one daily evaluation, market observations support forward labels, the full-market reference, research cohorts, and cash accounts under the same conventions. This limits semantic drift between financial research and benchmark reporting. Access and disclosure policies remain specific to each setting: a research interface's visible-date cutoff alone does not establish a sealed evaluation release. Likewise, common evaluation code supports consistency, while agreement of deployed entry points requires matched end-to-end run evidence.

The three input profiles are bound to actual releases before use. Raw-event processing, fifteen-second observations, and anonymous features introduce different preprocessing and information conditions; platform support alone does not establish that a dataset has been publicly released or used in a completed campaign. Publication receipts and coverage summaries identify the data underlying each leaderboard setting.

\subsection{Durable Accounting and Recoverable Progress}
\label{sec:arena_recovery}

Model calls, strategy registration, research evaluations, and final commitments carry persistent operation identities. The platform records budget reservations separately from observed consumption. Provider-reported counters settle model usage, and trusted execution measurements settle CPU usage when available. Unreported usage remains unknown. This separation keeps conservative admission estimates from appearing as measured costs.

A harness can checkpoint and yield while an evaluation runs, then resume from its saved research state. Strategy execution persists accepted daily outputs so that recoverable interruptions can continue the remaining dates. Completed model operations are recovered under their original identities. When a request's provider outcome or charge cannot be established, the operation remains unresolved pending reconciliation; it is not automatically replaced by another generation. Recovery thus preserves the candidate history and resource boundary of the original replicate.

The evaluator persists core financial and signal results before optional reference comparisons, reports, and detailed artifact publication. A failure in an optional step affects that evidence and any dependent component, while preserving the core result. Participant errors, infrastructure blockage, cancellation, and statistical unavailability remain distinct outcomes. These are properties of the execution design; campaign completion rates, recovery success, and runtime overhead are empirical quantities to be reported separately.

\subsection{Campaigns and Reproducibility Records}
\label{sec:arena_campaigns}

The campaign interface snapshots the declared participant source and public suite contract, persists submission intent, and submits through the shared Arena services. It collects evaluator-produced scorecards without recomputing financial results. Direct comparisons bind the service origin, suite, selected runtime, budget definition, and protocol. Different input scopes or financial conventions define separate groups.

A reproducibility record links the data publication, suite, participant source and prior information, actual model identity, model interactions, admitted trials, usage observations, final commitment, and final scorecard. All prescribed replicates remain visible in the campaign record, including failed, cancelled, blocked, and incomplete runs. Stored model replies support reconstruction of an observed trajectory; new calls to an external provider need not reproduce that trajectory exactly. The registration requirements appear in Appendix~\ref{app:registry}.
